\section{模型假设}
\begin{enumerate}
\item 每10分钟内功率按常数近似，负荷不可主动削减。该假设使连续能量平衡转化为有限维线性约束。
\item 储能充电和放电各按90\%效率计损，往返效率为81\%；不另计自放电与寿命衰减。该解释适用于题目给定的固定效率条件。
\item 微网不售电；无法被负荷或储能吸收的供给可弃置。预测规划中的弃光和执行中计划购电富余的弃电应区别理解，后者不获得退款。
\item 常规购电按已承诺电量结算，短缺额外以五倍当期价格采购。问题三各槽按最终有效调整量相对零点原计划结算一次，避免对同一槽重复计收偏差费用。
\item 问题一保留题面要求的日循环；问题二至四不人为固定每日实际末端电量，两日规划最远端仅受容量边界约束。日内调整规划的末端参考由当日两日计划内生给出，执行允许偏离该参考。
\item 波动电价主分析G将附件4给定价格作为规划窗口内已知输入；H作为仅有历史价格的扩展。G不是对未来实际负荷或光伏已知的假设。
\end{enumerate}
